\subsection{Adaptive soft-scaffold inpainting and beam search}
\label{sec:adaptive-soft-scaffold}

\subsubsection{Hard, soft, and free atom partitions}

The fixed-increment staged-growth experiments identified an addition of
four atoms per stage as the most effective granularity under the tested
protocol. The subsequent adaptive-flexibility experiment therefore held
the requested increment constant at \(n_{\mathrm{add}}=4\) and replaced
increment size with a stage-wise scaffold-rigidity action.

At stage \(s\), ligand atoms were partitioned into three disjoint sets.
The seven-heavy-atom warhead inherited from ligand A formed the hard set
\(M_{\mathrm{H}}\). Atoms generated and retained during earlier stages
formed the soft parent-scaffold set \(M_{\mathrm{S}}\). Newly generated
atoms formed the free set \(M_{\mathrm{U}}\). The complete supplied parent
scaffold was
\[
M_{\mathrm{F}}=M_{\mathrm{H}}+M_{\mathrm{S}},
\]
and \(M_{\mathrm{U}}=1-M_{\mathrm{F}}\).

The hard warhead remained fully inpainted throughout sampling. Soft
parent atoms retained their atom-feature channels but their coordinates
were only partially reintroduced, allowing previously generated geometry
to move without changing inherited atom identities.

\subsubsection{Soft-scaffold reinjection}

Let \(z_t^{\mathrm{known}}=(x_t^{\mathrm{known}},h_t^{\mathrm{known}})\)
denote the appropriately noised representation of the supplied parent
scaffold at reverse-diffusion step \(t\), and let
\(z_t^{\mathrm{model}}=(x_t^{\mathrm{model}},h_t^{\mathrm{model}})\)
denote the model state after the corresponding reverse transition. A
scalar rigidity coefficient \(\rho_s\in[0,1]\) controlled coordinate
reinjection for the soft parent atoms:
\[
x_t =
M_{\mathrm{H}}\odot x_t^{\mathrm{known}}
+
M_{\mathrm{S}}\odot
\left[
\rho_s x_t^{\mathrm{known}}
+
(1-\rho_s)x_t^{\mathrm{model}}
\right]
+
M_{\mathrm{U}}\odot x_t^{\mathrm{model}}.
\]
Atom-feature channels were combined as
\[
h_t =
(M_{\mathrm{H}}+M_{\mathrm{S}})\odot h_t^{\mathrm{known}}
+
M_{\mathrm{U}}\odot h_t^{\mathrm{model}}.
\]

Thus, \(\rho_s=1\) reproduced conventional hard inpainting of the
complete parent scaffold, whereas \(\rho_s=0\) released the coordinates
of the soft scaffold completely while preserving its atom types.
Intermediate values produced partial coordinate retention. The seven
warhead atoms were always assigned \(\rho=1\).

The \(\rho=1\) branch retained the original hard-inpainting execution
path. For \(\rho<1\), the per-sample ligand coordinate mean
\[
\mu_t =
\frac{1}{N_{\mathrm{lig}}}
\sum_{i=1}^{N_{\mathrm{lig}}}x_{t,i}
\]
was removed after every mixed reinjection, and the same translation was
applied to the corresponding pocket coordinates:
\[
x_{t,i}^{\mathrm{lig}}\leftarrow x_{t,i}^{\mathrm{lig}}-\mu_t,
\qquad
x_{t,j}^{\mathrm{pocket}}\leftarrow x_{t,j}^{\mathrm{pocket}}-\mu_t.
\]
This restored the zero-center-of-mass invariant required by the DiffSBDD
reverse process without altering internal ligand distances, internal
pocket distances, or ligand--pocket relative geometry.

\subsubsection{Order-safe scaffold propagation}

Parent identity was recovered using a one-to-one Hungarian assignment
constrained by atom type and coordinate distance. The seven hard atoms
were required to match their reference positions within 0.2~\AA. No
displacement ceiling was imposed on the soft parent atoms because their
movement was the experimental variable. Every retained child was then
renumbered as
\[
[\text{hard warhead}]
+
[\text{soft parent in parent order}]
+
[\text{new atoms}],
\]
both before molecular reconstruction and after RDKit processing. Parent
atom types were required to remain unchanged.

\subsubsection{Adaptive action search}

Each trajectory began from the bare seven-heavy-atom warhead. Stage 1
was sampled only with \(\rho_1=1\). At every later stage, the action set
was
\[
\mathcal{A}_{\rho}
=
\{1.00,0.75,0.50,0.25,0.00\}.
\]
For each live parent and each action, ten samples were generated with
\(n_{\mathrm{add}}=4\). All \(\rho\) values evaluated from the same parent
at the same stage used the same random seed, providing paired diffusion
noise across rigidity actions.

Sampling used the frozen model and the previously established settings:
shape-guidance strength \(\lambda=20\), Gaussian width \(\alpha=0.3\),
50 reverse-diffusion timesteps, five resampling iterations, and a
per-atom guidance clipping threshold of 1.0~\AA. No model parameters
were updated.

Greedy search retained one continuation per stage (\(k=1\)); beam search
retained up to three diverse continuations (\(k=3\)). Both variants used
identical generation, filtering, ranking, and stopping rules.

\subsubsection{Anchor-component filtering and parent retention}

Generated records were sanitized when possible and decomposed into
connected components. The anchor component was the component containing
all seven matched warhead atoms. A candidate was rejected when the
warhead could not be matched, exceeded the 0.2~\AA\ tolerance, or was
split across components.

For stages after the first, all inherited parent atoms were required to
remain in the anchor component with unchanged atom types. The anchor
child was additionally required to contain more heavy atoms than its
parent:
\[
H(C_{s+1})>H(C_s).
\]
Secondary components were discarded. Full-record connectivity was
retained as a separate observable and was not required for continuation.

\subsubsection{Shape gate, beam ranking, and diversity}

Shape was evaluated on the anchor component using RDKit Shape Tanimoto
distance \(d_{\mathrm{T}}\) and directional Shape Protrude distance
\(d_{\mathrm{P}}\), with \texttt{allowReordering=False}. Lower values
indicated greater similarity to ligand B.

For candidate child \(c\) and parent \(p\),
\[
g_{\mathrm{T}}
=
d_{\mathrm{T}}(p,B)-d_{\mathrm{T}}(c,B),
\qquad
g_{\mathrm{P}}
=
d_{\mathrm{P}}(p,B)-d_{\mathrm{P}}(c,B).
\]
After the initial stage, a candidate passed the trade-off gate when
\[
g_{\mathrm{T}}\geq-\varepsilon,
\qquad
g_{\mathrm{P}}\geq-\varepsilon,
\qquad
g_{\mathrm{T}}+g_{\mathrm{P}}>0,
\]
with \(\varepsilon=0.01\). The initial stage was filtered by structural
validity and anchor growth alone.

Passing candidates were ranked by absolute state quality,
\[
Q(c)
=
-\left[
d_{\mathrm{T}}(c,B)+d_{\mathrm{P}}(c,B)
\right].
\]
Selected candidates were required to differ by more than 0.75~\AA\ under
symmetric Chamfer distance calculated on newly grown heavy-atom
coordinates without structural alignment.

\subsubsection{Graduation and recorded outcomes}

The hard pair x0434-to-x2193 and the moderate pair x0874-to-x1093 had
target heavy-atom counts of 16 and 19, respectively. An anchor graduated
when its heavy-atom count lay within one atom of the corresponding
target. Anchors above the upper tolerance were rejected as oversize.

A component-graduated endpoint reached the target size in the anchor
component. A strictly connected endpoint additionally required the
complete generated record to contain one heavy-atom component.

Recorded quantities included the selected \(\rho\), cumulative \(\rho\)
schedule, anchor heavy-atom count, full-record connectivity, number of
heavy-atom components, warhead RMSD, complete-parent RMSD,
soft-scaffold RMSD, maximum parent-atom displacement, Shape Tanimoto
distance, Shape Protrude distance, local gain, cumulative gain, and
graduation status.

The final campaign comprised two ligand pairs, ten random seeds, and two
search variants, yielding 40 independent trajectories. Runs were
checkpointed after every completed stage and could terminate by
graduation, loss of all valid continuations, or the maximum of eight
stages.
